\documentclass[../readings.tex]{subfiles}

\begin{document}
\subsection{Koopman Theory, DMD, and SINDy}
\label{sub:koopman}

\textBox{%
Koopman analysis studies nonlinear dynamics through linear evolution of observables. DMD fits a finite-dimensional linear model from snapshots. SINDy fits sparse governing equations from a chosen function library.
}

\subsubsection{Koopman Operator Theory}

\addDef{Nonlinear Dynamics and Observables}{%
$\bx^{(k+1)} = \mathbf{F}(\bx^{(k)})$ on $\mathcal{M} \subseteq \fR^n$.
\begin{itemize}
    \item \textbf{Observable:}\enspace Scalar function $g: \mathcal{M} \to \fR$.
    \item \textbf{Koopman Operator $\mathcal{K}$:}\enspace Linear operator acting on observables: $(\mathcal{K}g)(\bx) = g(\mathbf{F}(\bx))$.
\end{itemize}
\textbf{Lifting:}\enspace $\mathcal{K}$ is linear even if $\mathbf{F}$ is nonlinear, at the cost of being infinite-dimensional.
}
\label{def:koopman-observables}

\addProp{Koopman Linearity}{%
The Koopman operator in Def~\ref{def:koopman-observables} is linear on observables:
\begin{align*}
    \mathcal{K}(\alpha g+\beta h)=\alpha\mathcal{K}g+\beta\mathcal{K}h.
\end{align*}
}
\label{prop:koopman-linearity}

\addPrf{%
For every state $\bx$,
\begin{align*}
    [\mathcal{K}(\alpha g+\beta h)](\bx)
    &=
    (\alpha g+\beta h)(\mathbf{F}(\bx))\\
    &=
    \alpha g(\mathbf{F}(\bx))+\beta h(\mathbf{F}(\bx))\\
    &=
    [\alpha\mathcal{K}g+\beta\mathcal{K}h](\bx).
\end{align*}
Since this holds for every $\bx$, the two observables are equal.
}

\addExmpl{%
\textbf{(Nonlinear map)}\enspace Let $x_{k+1}=x_k^2$. For observables $g_1(x)=x$ and $g_2(x)=x^2$,
\begin{align*}
    (\mathcal{K}g_1)(x)=x^2=g_2(x),
    \qquad
    (\mathcal{K}g_2)(x)=x^4.
\end{align*}
The state map is nonlinear, but the operator $g\mapsto g\circ F$ is linear on the space of observables.
}

\addDef{Koopman Eigenfunctions}{%
Observable $\varphi$ satisfying $\mathcal{K}\varphi = \mu \varphi$.
\begin{itemize}
    \item \textbf{Temporal Evolution:}\enspace $\varphi(\bx^{(k)}) = \mu^k \varphi(\bx^{(0)})$.
    \item \textbf{Significance:}\enspace Provides a global linear coordinate system for nonlinear dynamics.
\end{itemize}
}
\label{def:koopman-eigenfunctions}

\subsubsection{Dynamic Mode Decomposition (DMD)}

\addDef{DMD Algorithm}{%
Finds the best linear fit $\bY \approx \bA\bX$ for snapshot matrices:
\begin{align*}
    \bX = [\bx^{(0)} \dots \bx^{(m-1)}], \quad \bY = [\bx^{(1)} \dots \bx^{(m)}].
\end{align*}
\begin{enumerate}
    \item Compute truncated SVD: $\bX \approx \bU_r \bsigma_r \bV_r^T$.
    \item Construct $r \times r$ reduced operator: $\tilde{\bA} = \bU_r^T \bY \bV_r \bsigma_r^{-1}$.
    \item DMD eigenvalues $\{\lambda_i\}$ and modes $\{\boldsymbol{\phi}_i\}$ are recovered from $\tilde{\bA}$.
\end{enumerate}
}
\label{def:dmd-algorithm}

\addThm{DMD Least Squares Model}{%
The best-fit linear map in Frobenius norm is
\begin{align*}
    \bA_*=\bY\bX^\dagger.
\end{align*}
The reduced operator in Def~\ref{def:dmd-algorithm} is the projection of this map onto the rank-$r$ left singular subspace of $\bX$.
}
\label{thm:dmd-least-squares}

\addExcr{%
\begin{enumerate}
    \item Use the least squares normal equations to derive $\bA_*=\bY\bX^\dagger$.
    \item Substitute the truncated SVD $\bX\approx\bU_r\bsigma_r\bV_r^T$.
    \item Show that $\bU_r^T\bA_*\bU_r=\bU_r^T\bY\bV_r\bsigma_r^{-1}$.
    \item Interpret this as the small operator $\tilde{\bA}$ in Def~\ref{def:dmd-algorithm}.
\end{enumerate}
}
\label{ex:dmd-least-squares}

\addRemark{%
\textbf{Interpretation:}\enspace DMD decomposes data into spatiotemporal modes $\boldsymbol{\phi}_i$ with associated frequencies and growth rates. It is the data-driven approximation of the Koopman operator using the identity dictionary.
}

\subsubsection{SINDy: Sparse Identification}

\addDef{Sparse Regression for Dynamics}{%
Identifies governing equations $\dot{\bx} = \mathbf{\Xi}^T \boldsymbol{\Theta}(\bx)$ where $\boldsymbol{\Theta}$ is a library of candidate functions (monomials, sin, etc.).
\begin{itemize}
    \item \textbf{STLS Algorithm:}\enspace Sequential Thresholded Least Squares. Iteratively solves least squares and thresholds small coefficients to promote sparsity.
    \item \textbf{Outcome:}\enspace Parsimonious, interpretable models (e.g., recovering Lorenz equations).
\end{itemize}
}
\label{def:sindy}

\addRemark{%
\textbf{(DMD vs.\ SINDy)}\enspace DMD fits a linear time-advance map for snapshot data. SINDy fits a sparse differential equation in a chosen library. DMD returns modes and growth rates; SINDy returns equation coefficients.
}

\addRemark{%
\textbf{(Noisy derivatives)}\enspace SINDy is sensitive to derivative estimates. Smoothing, weak-form regression, or integral formulations are often needed when data are noisy.
}

\subsubsection{Exercises}

\addExcr{%
\begin{enumerate}
    \item Prove Prop~\ref{prop:koopman-linearity} without looking at the proof.
    \item Implement DMD for a 2D damped oscillator. Verify the identified eigenvalues.
    \item Build a polynomial library $\boldsymbol{\Theta}$ for $\bx \in \fR^2$ up to degree 2.
    \item Use SINDy (via STLS) to identify the Lotka-Volterra predator-prey model from noisy data.
    \item Compare DMD and SINDy on the same two-dimensional linear system. Which quantities should agree?
\end{enumerate}
}

\end{document}
